\documentclass[a4paper, serif, 10pt]{moderncv}

%% moderncv
% color and style
\moderncvcolor{black}
\moderncvstyle{banking}

% renew moderncv command to adapt for CJK colon
\renewcommand*{\cvitem}[3][.25em]{%
  \ifstrempty{#2}{}{\hintstyle{#2}：}{#3}%
  \par\addvspace{#1}}

\renewcommand*{\cvitemwithcomment}[4][.25em]{%
  \savebox{\cvitemwithcommentmainbox}{\ifstrempty{#2}{}{\hintstyle{#2}：}#3}%
  \setlength{\cvitemwithcommentmainlength}{\widthof{\usebox{\cvitemwithcommentmainbox}}}%
  \setlength{\cvitemwithcommentcommentlength}{\maincolumnwidth-\separatorcolumnwidth-\cvitemwithcommentmainlength}%
  \begin{minipage}[t]{\cvitemwithcommentmainlength}\usebox{\cvitemwithcommentmainbox}\end{minipage}%
  \hfill% fill of \separatorcolumnwidth
  \begin{minipage}[t]{\cvitemwithcommentcommentlength}\raggedleft\small\itshape#4\end{minipage}%
  \par\addvspace{#1}}

%% page layout/margins
\usepackage[top=2.5cm, bottom=2.5cm, left=1.5cm, right=1.5cm]{geometry}

\nopagenumbers{}

%% Babel config for English language
\usepackage[english]{babel}

%% fontspec
\usepackage{fontspec}

\IfFontExistsTF{Linux Libertine}{
  \setmainfont[Ligatures={TeX, Common}, Numbers=Lining, ItalicFont=Linux Libertine]{Linux Libertine}
}{}
\IfFontExistsTF{Linux Libertine O}{
  \setmainfont[Ligatures={TeX, Common}, Numbers=Lining, ItalicFont=Linux Libertine O]{Linux Libertine O}
}{}

%% CTeX
% CJK support, used to show CJK characters in the resume
%
% - heading=false: disable ctex heading
% - punct=kaiming: use kaiming punctuations styles for CJK
% - scheme=plain: use plain scheme, do not override \normalsize font size
% - space=auto: space settings for CJK characters
%
% CJK fonts are resolved per locale language, e.g. Noto Serif/Sans CJK JP
% for Japanese resumes, with fallback to the next region if not installed
%
% ref:
% - http://ctan.mirrorcatalogs.com/language/chinese/ctex/ctex.pdf
\usepackage[UTF8, heading=false, punct=kaiming, scheme=plain, space=auto]{ctex}

\IfFontExistsTF{Noto Serif CJK SC}{
  \setCJKmainfont{Noto Serif CJK SC}
}{}
\IfFontExistsTF{Noto Sans CJK SC}{
  \setCJKsansfont{Noto Sans CJK SC}
}{}

%% line spacing
\usepackage{setspace}
\setstretch{1.125}

%% URL styling - use normal text instead of monospace
\urlstyle{same}

%% Patch \httplink and \httpslink to handle full URLs
% The original moderncv macros always prepend http:// or https:// to URLs.
% This causes issues when the URL already has a protocol (e.g., https://example.com)
% resulting in malformed URLs like https://https://example.com.
% This patch checks if the URL already contains "://" and uses it directly if so.
% Note: We use \str_set:Nx (x-expansion) to expand macros like \@homepage first.
\ExplSyntaxOn
\str_new:N \g_yamlresume_url_str

\RenewDocumentCommand{\httpslink}{O{}m}{
  \str_set:Nx \g_yamlresume_url_str {#2}
  \str_if_in:NnTF \g_yamlresume_url_str {://}
    { \url{#2} }
    { \url{https://#2} }
}

\RenewDocumentCommand{\httplink}{O{}m}{
  \str_set:Nx \g_yamlresume_url_str {#2}
  \str_if_in:NnTF \g_yamlresume_url_str {://}
    { \url{#2} }
    { \url{http://#2} }
}
\ExplSyntaxOff

\name{陈思远}{}
\title{资深业务分析师 | 数字化转型与数据驱动决策}
\phone[mobile]{+86 21 6888 0136}
\email{siyuan.chen@example.com}
\homepage{https://www.siyuanchen-analytics.com}

\address{中国，上海市，上海，浦东新区世纪大道 100 号，200120}{}{}

\extrainfo{{\small \faLinkedin}\ \href{https://www.linkedin.com/in/siyuan-chen-ba}{@siyuan-chen-ba} {} {} {} • {} {} {} 
{\small \faMedium}\ \href{https://medium.com/@siyuan.chen}{@siyuan.chen} {} {} {} • {} {} {} 
{\small \faZhihu}\ \href{https://www.zhihu.com/people/siyuan-chen-ba}{@陈思远}}

\begin{document}

\maketitle

\section{简介}

\cvline{}{\begin{itemize}
\item 9 年金融科技与供应链领域业务分析经验，擅长需求调研、流程建模、数据洞察与跨部门协同。
\item 主导 10+ 个数字化项目，覆盖业务蓝图、需求规格、UAT 测试与上线复盘，推动流程效率提升 25\% 以上。
\item 熟练使用 SQL、Tableau、Power BI、Jira 与 Confluence，能够将复杂业务问题转化为可执行方案。
\item 具备优秀的利益相关方管理与沟通能力，善于在业务、产品、技术与运营团队之间建立高效协作。
\end{itemize}}
\section{教育背景}

\cventry{2012年9月–2016年6月}
        {学士，信息管理与信息系统，成绩：3.8/4.0}
        {上海财经大学}
        {\url{https://www.sufe.edu.cn}}
        {}
        {\begin{itemize}
\item 系统学习信息系统规划、数据库设计与业务流程优化
\item 获校级优秀学生奖学金，毕业设计聚焦企业数据治理与流程分析
\end{itemize}
\textbf{课程}：管理信息系统、数据库原理与应用、运筹学、统计学、业务流程管理、项目管理、数据分析基础、经济学原理}
\section{工作经历}

\cventry{2021年7月至今}
        {高级业务分析师}
        {星瀚数字科技（上海）有限公司}
        {\url{https://www.xinghan-digital.com}}
        {}
        {\begin{itemize}
\item 负责金融科技产品线的业务分析与需求管理，覆盖客户尽调、风控审批与运营监控等核心流程
\item 主导 6 个跨部门数字化项目，输出业务蓝图、需求规格说明书、流程模型与 UAT 方案
\item 建立需求优先级评估机制，协调产品、研发、风控与运营团队，按期交付率提升至 95\%
\item 设计经营分析看板与关键指标口径，支持管理层进行月度经营复盘与资源调配
\item 组织用户访谈与工作坊，沉淀需求模板与业务术语库，缩短新成员上手周期 30\%
\end{itemize}
\textbf{关键字}：需求管理、流程建模、利益相关方沟通、数据看板、金融科技}

\cventry{2018年4月–2021年6月}
        {业务分析师}
        {云帆供应链管理有限公司}
        {\url{https://www.yunfan-supply.com}}
        {}
        {\begin{itemize}
\item 深入采购、仓储与物流业务，梳理订单履约全链路流程并识别效率瓶颈
\item 完成 20+ 份业务需求文档与流程现状图，推动仓储作业标准化和异常处理机制落地
\item 与数据团队合作搭建库存周转、履约时效与成本分析报表，辅助运营策略调整
\item 参与 ERP 系统升级项目，负责需求澄清、测试用例评审与上线培训
\item 通过流程优化将订单异常处理时长缩短 18\%，库存盘点差异率降低 12\%
\end{itemize}
\textbf{关键字}：供应链、流程优化、ERP、需求文档、数据分析}

\cventry{2016年7月–2018年3月}
        {数据分析师}
        {睿联市场咨询有限公司}
        {\url{https://www.ruilian-consulting.com}}
        {}
        {\begin{itemize}
\item 为零售与消费品客户提供市场研究、用户洞察与经营分析支持
\item 使用 SQL 与 Excel 清洗和分析调研数据，撰写行业分析与客户洞察报告
\item 参与 15+ 个咨询项目，负责访谈纪要、问卷设计、数据验证与汇报材料制作
\item 与项目经理共同向客户汇报关键发现，提出可落地的业务改进建议
\end{itemize}
\textbf{关键字}：市场研究、用户洞察、SQL、报告撰写、零售行业}
\section{语言}

\cvline{中文}{母语或双语水平 \hfill \textbf{关键字}：普通话二级甲等、商务写作}
\cvline{英语}{完全专业水平 \hfill \textbf{关键字}：CET-6、商务沟通、英文报告撰写}
\cvline{日语}{初级水平 \hfill \textbf{关键字}：JLPT N4、基础商务会话}
\section{专业技能}

\cvline{业务分析}{专家 \hfill \textbf{关键字}：需求调研与访谈、业务流程建模、需求规格说明书、用户故事与用例、UAT 测试管理、利益相关方管理}
\cvline{数据分析}{高级 \hfill \textbf{关键字}：SQL、Tableau、Power BI、Excel 高级分析、指标体系设计、A/B 测试基础}
\cvline{产品与项目}{高级 \hfill \textbf{关键字}：Jira、Confluence、Visio、敏捷开发、项目计划与跟踪、风险与变更管理}
\cvline{行业知识}{高级 \hfill \textbf{关键字}：金融科技、供应链管理、零售与消费品、数字化转型、数据治理}
\section{荣誉}

\cventry{2023年12月}
        {星瀚数字科技（上海）有限公司}
        {年度优秀员工}
        {}
        {}
        {表彰在金融科技业务分析项目中的突出贡献，主导需求管理机制建设并推动跨部门协作效率提升。}

\cventry{2020年11月}
        {云帆供应链管理有限公司}
        {业务分析卓越奖}
        {}
        {}
        {因供应链流程优化项目表现优异获奖，项目实现订单异常处理时长缩短 18\%。}
\section{证书}

\cventry{2022年9月}
        {国际商业分析协会（IIBA）}
        {国际商业分析专业人士认证（CBAP）}
        {\url{https://www.iiba.org/certification/cbap}}
        {}
        {}

\cventry{2021年6月}
        {美国项目管理协会（PMI）}
        {项目管理专业人士认证（PMP）}
        {\url{https://www.pmi.org/certifications/project-management-pmp}}
        {}
        {}

\cventry{2020年8月}
        {谷歌}
        {谷歌数据分析专业证书}
        {\url{https://grow.google/certificates/data-analytics/}}
        {}
        {}
\section{出版刊物}

\cventry{2023年5月}
        {数据驱动的业务流程优化实践}
        {企业管理评论}
        {\url{https://www.qyglpl.com/article/2023/data-driven-process-optimization}}
        {}
        {\begin{itemize}
\item 结合供应链与金融科技项目案例，总结从业务问题识别到指标落地的分析方法
\item 提出以流程数据、用户反馈与经营指标相结合的需求优先级评估框架
\end{itemize}}
\section{引荐}

\cventry{\emaillink[min.zhou@example.com]{min.zhou@example.com}}
        {星瀚数字科技产品总监}
        {周敏}
        {+86 21 6888 0142}
        {}
        {陈思远具备扎实的业务分析能力与出色的跨部门沟通能力，能够快速理解复杂业务并推动方案落地。}

\cventry{\emaillink[hao.li@example.com]{hao.li@example.com}}
        {云帆供应链运营负责人}
        {李昊}
        {+86 21 5566 0189}
        {}
        {在供应链流程优化项目中展现出严谨的分析思维和较强的推动力，是团队值得信赖的业务伙伴。}
\section{项目}

\cventry{2019年3月–2020年10月}
        {面向采购、仓储与物流团队的履约全链路可视化与异常预警平台}
        {供应链履约可视化平台}
        {\url{https://www.yunfan-supply.com/projects/visibility-platform}}
        {}
        {\begin{itemize}
\item 负责业务需求调研、流程现状梳理与需求规格编写
\item 设计订单履约、库存周转与异常处理等核心指标口径
\item 协调研发与运营完成 UAT 测试及上线培训
\item 上线后订单异常平均处理时长缩短 18\%
\end{itemize}
\textbf{关键字}：供应链、可视化、异常预警、需求分析}

\cventry{2022年2月–2023年8月}
        {覆盖客户准入、风险评级与审批协同的尽调流程数字化项目}
        {客户尽调流程数字化}
        {\url{https://www.xinghan-digital.com/projects/kyc-digital}}
        {}
        {\begin{itemize}
\item 访谈风控、合规与客户经理，梳理尽调流程痛点与合规要求
\item 输出业务蓝图、用户故事、流程模型与验收标准
\item 设计审批时效与退回率等运营指标，支持持续优化
\item 项目上线后平均审批周期缩短 22\%
\end{itemize}
\textbf{关键字}：金融科技、客户尽调、流程数字化、合规}

\cventry{2017年3月–2017年12月}
        {基于调研与销售数据的零售客户分群与经营策略分析项目}
        {零售客户洞察分析}
        {\url{https://www.ruilian-consulting.com/projects/retail-insight}}
        {}
        {\begin{itemize}
\item 设计问卷与访谈提纲，完成多城市消费者调研数据清洗
\item 使用 SQL 与 Excel 进行客户分群、交叉分析与趋势洞察
\item 撰写客户洞察报告并向客户管理层汇报策略建议
\end{itemize}
\textbf{关键字}：零售、客户洞察、市场研究、数据分析}
\section{兴趣爱好}

\cvline{行业研究}{数字化转型、商业分析、组织效能}
\cvline{运动}{马拉松、游泳、徒步}
\cvline{创作}{商业写作、摄影、播客}
\section{志愿服务}

\cventry{2020年9月–2024年6月}
        {公益导师}
        {上海青年商业分析社区}
        {\url{https://www.shyba-community.org}}
        {}
        {\begin{itemize}
\item 为高校学生与初入职场者提供商业分析职业发展与案例辅导
\item 参与组织需求分析工作坊与行业分享活动
\item 协助社区沉淀案例库与导师制流程
\end{itemize}}
    
\end{document}